% Experimental margin engine. Candidate geometry only; not a public profile yet.
\documentclass[
  11pt,
  a4paper,
  twoside,
  open=any,
  DIV=11
]{scrbook}

\usepackage[UTF8,fontset=none,heading=false,scheme=chinese]{ctex}
\usepackage{ipara-core}
\usepackage{scrlayer-scrpage}
\usepackage{marginnote}

% Candidate margin geometry. These numbers are experimental, not contract values.
\KOMAoptions{areasetadvanced=true,headinclude=true,footinclude=false,headheight=18pt,footheight=18pt}
\areaset{116mm}{220mm}
\setlength{\marginparsep}{5mm}
\setlength{\marginparwidth}{39mm}
\setlength{\marginparpush}{7pt}

\setkomafont{chapter}{\normalfont\huge\bfseries\sffamily\color{iparaDeep}}
\setkomafont{section}{\normalfont\Large\bfseries\sffamily\color{iparaDeep}}
\setkomafont{subsection}{\normalfont\large\bfseries\sffamily\color{iparaDeep}}
\setkomafont{pageheadfoot}{\normalfont}
\setkomafont{pagehead}{\normalfont}
\setkomafont{pagefoot}{\normalfont}
\RedeclareSectionCommand[beforeskip=1.15em,afterskip=0.55em]{section}
\renewcommand*{\chapterformat}{第\,\thechapter\,章\autodot\enskip}

\clearpairofpagestyles
\ihead{\small\sffamily I.P.A.R.A · Margin Engine}
\ohead{\small\sffamily\headmark}
\cfoot{\small\pagemark}
\automark[section]{chapter}
\pagestyle{scrheadings}
\renewcommand*{\chapterpagestyle}{plain.scrheadings}

\raggedbottom
\setlength{\parindent}{2em}
\setlength{\parskip}{0.28em plus 0.08em minus 0.04em}

\newcommand{\semanticmargin}[2]{%
  \marginnote{%
    \raggedright\footnotesize\sffamily
    \textbf{#1}\par
    \vspace{0.2em}%
    \normalfont\footnotesize #2%
  }%
}

% KOMA addmargin* automatically mirrors the extended side on odd/even pages.
% V1 widecontent is intentionally a local block; crossing pages should warn.
\NewDocumentEnvironment{widecontent}{+b}{%
  \par\addvspace{0.55em}%
  \begin{addmargin*}[0pt]{-\dimexpr\marginparsep+\marginparwidth\relax}%
    #1%
  \end{addmargin*}%
  \par\addvspace{0.55em}%
}{}

\hypersetup{
  unicode=true,
  colorlinks=true,
  linkcolor=iparaDeep,
  urlcolor=iparaAccent,
  pdftitle={I.P.A.R.A Margin Engine Specimen}
}

\makeatletter
\AtBeginDocument{%
  \typeout{IPARA-MARGIN-TEXTWIDTH=\the\textwidth}%
  \typeout{IPARA-MARGIN-TEXTHEIGHT=\the\textheight}%
  \typeout{IPARA-MARGIN-ODDSIDEMARGIN=\the\oddsidemargin}%
  \typeout{IPARA-MARGIN-EVENSIDEMARGIN=\the\evensidemargin}%
  \typeout{IPARA-MARGIN-MARGINPARWIDTH=\the\marginparwidth}%
  \typeout{IPARA-MARGIN-MARGINPARSEP=\the\marginparsep}%
}
\makeatother

\begin{document}

\frontmatter
\tableofcontents

\mainmatter
\chapter{MainText 与 Margin 的职责}

\section{MainText 必须逻辑闭合}

\semanticmargin{Owner}{MainText owns the reasoning chain. Margin only carries navigation, reminders and local references.}

设线性映射 $T\colon V\to W$。选基以后得到矩阵 $A$。如果读者完全遮住本页边栏，这一段仍然必须完整说明：矩阵是映射的表示，而不是映射本身；合法换基改变表示，但不改变抽象映射。

\semanticmargin{Boundary}{表示变化 $\neq$ 对象变化。边栏可以提醒这个边界，但第一次完整定义必须留在正文。}

第二段继续故意写成无需边栏才能理解的形式。Margin 的价值是降低检索和回忆成本，不是把正文切碎成“必须左右来回看”的网页 UI。

\begin{mentalmodel}{Mental Model · Margin 是辅助通道}
MainText 构成可独立阅读的知识主链；semantic margin 只承载不影响下一步推理成立的补充信息。
\end{mentalmodel}

\section{Wide content 不应该靠缩字}

下面的表格故意超过 112mm 主阅读栏。它应临时消费主栏与外侧 margin 的总包络，而不是缩成很小的字体。

\begin{widecontent}
\begin{iparatable}{L{0.18\linewidth}YYY}
\toprule
\textbf{State} & \textbf{MainText} & \textbf{Margin} & \textbf{Large object}\\
\midrule
main & 完整推理主链 & 导航、提醒、小型引用 & 服从主栏宽度\\
widecontent & 仍是当前正文对象 & 暂时让出空间 & 使用主栏 + 外侧 margin 包络\\
fullwidth & 特殊全景对象 & 暂停 ordinary margin & 后续单独验证，不在本实验实现\\
\bottomrule
\end{iparatable}
\end{widecontent}

\newpage
\section{偶数页镜像检查}

这一页由 \texttt{twoside} 强制进入另一侧页面。`addmargin*` 必须把 `widecontent` 自动扩到新的外侧，而不是继续向同一个物理方向伸出。

\semanticmargin{Even-page check}{如果页码方向翻转，semantic margin 也必须位于新的 outer margin。}

\begin{widecontent}
\begin{iparatable}{L{0.18\linewidth}YYY}
\toprule
\textbf{Odd/Even} & \textbf{Inner side} & \textbf{Outer side} & \textbf{Expected behavior}\\
\midrule
Odd & 靠装订侧 & semantic margin & widecontent 向 outer side 扩展\\
Even & 靠装订侧镜像 & semantic margin 镜像 & widecontent 随 outer side 镜像\\
\bottomrule
\end{iparatable}
\end{widecontent}

\chapter{分页边界}

\section{Wide block 不跨页}

V1 把 `widecontent` 定义为局部状态。KOMA 的 `addmargin*` 会在 block 跨页时发出潜在错误边距警告，因此大型多页表格不能偷偷塞进这个环境；应单独设计 full-width / long-table policy。

\semanticmargin{Stop Boundary}{`widecontent` 是局部逃逸，不是多页版式系统。}

\end{document}
